% =============================================================
% style.tex - Farben, Fußnoten, Code-Highlighting, Emojis, Hyperref
% Ersetzt: Farben_definieren-settings.tex + Fussnoten-settings.tex +
%          Listings-settings.tex + ASCII-Art-settings.tex +
%          Emoji-settings.tex + Hyperref-settings.tex
% =============================================================

% =============================================================
% 1. FARBEN
% =============================================================

% Hauptpalette
\definecolor{primary}{RGB}{0, 102, 153}
\definecolor{secondary}{RGB}{102, 102, 102}
\definecolor{accent}{RGB}{204, 51, 0}

% Grautöne
\definecolor{darkgray}{RGB}{64, 64, 64}
\definecolor{lightgray}{RGB}{192, 192, 192}
\definecolor{verylightgray}{RGB}{240, 240, 240}

% Semantische Farben
\definecolor{success}{RGB}{40, 167, 69}
\definecolor{warning}{RGB}{255, 193, 7}
\definecolor{danger}{RGB}{220, 53, 69}
\definecolor{info}{RGB}{23, 162, 184}

% Code-Farben
\definecolor{codebackground}{RGB}{248, 248, 248}
\definecolor{codebg}{RGB}{248, 248, 248}
\definecolor{codeborder}{RGB}{224, 224, 224}

% UniBw-Farben
\definecolor{unibwblue}{RGB}{0, 51, 102}
\definecolor{unibwgray}{RGB}{102, 102, 102}

% Digital vs. Druck (Schalter \ifisprint in 1Main.tex)
\ifisprint
  \definecolor{printblack}{cmyk}{0,0,0,0.85}
  \definecolor{printblue}{cmyk}{1,0.5,0,0}
  \colorlet{textcolor}{printblack}
  \colorlet{linkcolor}{printblue}
  \colorlet{headingcolor}{printblack}
\else
  \colorlet{textcolor}{black}
  \colorlet{linkcolor}{primary}
  \colorlet{headingcolor}{primary}
\fi

% Utility-Befehle
\newcommand{\textprimary}[1]{\textcolor{primary}{#1}}
\newcommand{\textsecondary}[1]{\textcolor{secondary}{#1}}
\newcommand{\textaccent}[1]{\textcolor{accent}{#1}}
\newcommand{\highlight}[1]{\colorbox{verylightgray}{#1}}

% TikZ-Farbstile
\tikzset{
  primary/.style={color=primary, fill=primary!10},
  secondary/.style={color=secondary, fill=secondary!10},
  accent/.style={color=accent, fill=accent!10},
  success/.style={color=success, fill=success!10},
  danger/.style={color=danger, fill=danger!10}
}

% =============================================================
% 2. FUßNOTEN
% =============================================================

\usepackage[bottom, hang, stable, multiple, norule]{footmisc}

\setlength{\footnotemargin}{12pt}
\setlength{\footnotesep}{8pt}
\setlength{\skip\footins}{12pt plus 4pt minus 2pt}

\renewcommand{\thefootnote}{\textcolor{secondary}{\arabic{footnote}}}
\renewcommand{\footnoterule}{%
  \kern-3pt
  {\color{primary}\hrule width 0.4\columnwidth height 0.5pt}
  \kern 2.6pt
}
\addtokomafont{footnote}{\color{darkgray}}
\interfootnotelinepenalty=10000

\usepackage{footnotebackref}
\usepackage{tablefootnote}

% Fußnoten-Hilfsbefehle
\newcommand{\sourcefootnote}[1]{\footnote{\textit{Quelle:} #1}}
\newcommand{\authornote}[1]{\footnote{\textit{Anm. d. Verf.:} #1}}
\newcommand{\transnote}[1]{\footnote{\textit{Übersetzung d. Verf.:} #1}}

% =============================================================
% 3. CODE-HIGHLIGHTING (MINTED)
% =============================================================

\PassOptionsToPackage{unicode=true, breaklinks=true, colorlinks=true}{hyperref}

\usepackage[
  cache=true,
  cachedir=build/minted-cache,
  frozencache=false
]{minted}

\setminted{
  bgcolor=codebg,
  frame=single,
  framesep=10pt,
  rulecolor=codeborder,
  fontsize=\small,
  fontfamily=tt,
  breaklines=true,
  breakanywhere=true,
  autogobble=true,
  linenos=false,
  style=default,
  tabsize=2
}

% Sprachspezifische Configs (linenos für Programmiersprachen)
\setminted[python]{python3=true, linenos=true, numbersep=8pt}
\setminted[java]{linenos=true, numbersep=8pt}
\setminted[scala]{linenos=true, numbersep=8pt}
\setminted[javascript]{linenos=true, numbersep=8pt}
\setminted[rust]{linenos=true, numbersep=8pt}
\setminted[json]{linenos=false}
\setminted[yaml]{linenos=false}

% Inline-Code und Datei-Imports
\newcommand{\inlinecode}[2][text]{\mintinline{#1}{#2}}
\newcommand{\inputcode}[3][]{\inputminted[#1]{#2}{#3}}

% Terminal-Umgebung für ASCII-Art / Konsolenausgaben
\IfFontExistsTF{DejaVu Sans Mono}{
  \newfontfamily\terminalfont{DejaVu Sans Mono}[Scale=0.85]
}{
  \newcommand{\terminalfont}{\ttfamily}
}
\newminted[terminal]{text}{
  formatcom=\terminalfont\linespread{0.95},
  fontsize=\footnotesize,
  breaklines=true,
  frame=single,
  framesep=10pt,
  rulecolor=codeborder,
  bgcolor=codebg
}

% Float-Umgebung für Code mit Caption
\usepackage{newfloat}
\DeclareFloatingEnvironment[
  fileext=lol,
  listname={Codeverzeichnis},
  name=Code,
  placement=tbp,
]{codefloat}

% =============================================================
% 4. EMOJIS
% =============================================================

\newif\ifEmojiFontFound
\EmojiFontFoundfalse

\IfFontExistsTF{Noto Color Emoji}{
  \newfontfamily{\emojifont}{Noto Color Emoji}[Renderer=HarfBuzz, Scale=MatchUppercase]
  \EmojiFontFoundtrue
}{
  \IfFontExistsTF{Apple Color Emoji}{
    \newfontfamily{\emojifont}{Apple Color Emoji}[Renderer=HarfBuzz, Scale=MatchUppercase]
    \EmojiFontFoundtrue
  }{
    \let\emojifont\rmfamily
  }
}

\newcommand{\emoji}[1]{{\emojifont #1}}

% Status-Symbole (mit TikZ-Fallback)
\makeatletter
\@ifpackageloaded{tikz}{
  \newcommand{\drawSuccess}{%
    \begin{tikzpicture}[baseline=-0.5ex, scale=0.25]
      \fill[success] (0,0) circle (1.2);
      \draw[white, line width=0.5mm] (-0.5,0) -- (-0.2,-0.4) -- (0.6,0.5);
    \end{tikzpicture}%
  }
  \newcommand{\drawFailure}{%
    \begin{tikzpicture}[baseline=-0.5ex, scale=0.25]
      \fill[danger] (0,0) circle (1.2);
      \draw[white, line width=0.5mm] (-0.4,-0.4) -- (0.4,0.4);
      \draw[white, line width=0.5mm] (-0.4,0.4) -- (0.4,-0.4);
    \end{tikzpicture}%
  }
}{
  \newcommand{\drawSuccess}{\textcolor{success}{Ok}}
  \newcommand{\drawFailure}{\textcolor{danger}{X}}
}
\makeatother

\ifEmojiFontFound
  \providecommand{\successmark}{\textcolor{success}{\emoji{✅}}}
  \providecommand{\failuremark}{\textcolor{danger}{\emoji{❌}}}
\else
  \providecommand{\successmark}{\drawSuccess}
  \providecommand{\failuremark}{\drawFailure}
\fi

% Häufige Text-Hilfsbefehle
\newcommand{\important}[1]{\textbf{\textcolor{warning}{⚠ Wichtig:}} #1}
\newcommand{\note}[1]{\textbf{\textcolor{info}{Hinweis:}} #1}
\newcommand{\tip}[1]{\textbf{\textcolor{primary}{Tipp:}} #1}

% =============================================================
% 5. HYPERREF & PDF-METADATEN
% =============================================================

% hyperref wird von minted geladen — nur noch konfigurieren
\hypersetup{
  unicode=true,
  breaklinks=true,
  colorlinks=true,
  linkcolor=blue!70!black,
  citecolor=green!60!black,
  urlcolor=blue!80!black,
  filecolor=purple!60!black,
  anchorcolor=black,
  bookmarks=true,
  bookmarksopen=true,
  bookmarksnumbered=true,
  bookmarksopenlevel=2,
  pdfstartview=FitH,
  pdfdisplaydoctitle=true,
  pdfnewwindow=true,
  pdfpagelayout=OneColumn,
  pdfpagemode=UseOutlines,
  hidelinks=false,
  pdflang={de-DE},
  pdfcreator={LuaLaTeX},
  pdfproducer={LaTeX with hyperref and LuaTeX}
}

% PDF-Metadaten aus metadata.tex – erst nach % =============================================================
% metadata.tex – Alle persönlichen Angaben HIER eintragen!
% =============================================================

% --- Autor ---
\newcommand{\myname}{Vorname Nachname}
\newcommand{\matricle}{Matrikelnummer}


% --- Arbeit ---
\newcommand{\mytitle}{Titel der Arbeit}
\newcommand{\mytype}{Bachelorarbeit}           % z.B. Bachelorarbeit / Masterarbeit / Seminararbeit

% --- Institution ---
\newcommand{\myuniversity}{Universität Regensburg}
\newcommand{\myfaculty}{Fakultät für Wirtschaftswissenschaften}
\newcommand{\myinstitute}{Institut für Angewandte Informatik}

% --- Betreuer ---
\newcommand{\reviewerone}{Prof. Dr. Mustermann}
\newcommand{\reviewertwo}{Prof. Dr. Mustermann}
\newcommand{\advisor}{Dr. Musterfrau}
        
% --- Datum ---
\newcommand{\timeend}{\today}

% --- LaTeX-Metadaten (automatisch befüllt, nichts ändern) ---
\author{\myname}
\title{\mytitle}
 auswerten
\AtBeginDocument{%
  \hypersetup{%
    pdfauthor={\myname},%
    pdftitle={\mytitle},%
    pdfsubject={\mytype}%
  }%
}


\def\UrlBreaks{%
    \do\-\do\_\do\/\do\*\do\~\do\'\do\"
    \do\.\do\:\do\;\do\,\do\?\do\!\do\=
    \do\+\do\&\do\|\do\@\do\#
}

\usepackage{hyperxmp}
\urlstyle{same}

\usepackage{bookmark}
\bookmarksetup{numbered, open, openlevel=2}

% Deutsche Namen für \autoref – hyperref bringt mit polyglossia nur englische mit
\def\chapterautorefname{Kapitel}
\def\sectionautorefname{Abschnitt}
\def\subsectionautorefname{Abschnitt}
\def\subsubsectionautorefname{Abschnitt}
\def\paragraphautorefname{Absatz}
\def\appendixautorefname{Anhang}
\def\figureautorefname{Abbildung}
\def\tableautorefname{Tabelle}
\def\equationautorefname{Gleichung}
\def\footnoteautorefname{Fußnote}
\def\itemautorefname{Punkt}
\def\definitionautorefname{Definition}
\def\listingautorefname{Listing}

